\paragraph{T/B $\dn$：压缩。}
单式T/B等于token数除以原串字节数。表\ref{tab:example}的公式A有11字节，软AST-BPE与AST-BPE均分成5个token：
\[\mathrm{T/B}_A=5/11\approx0.4545.\]
Qwen2.5使用7个token，对应$7/11\approx0.6364$。实验表先计算每个来源组的总token数／总字节数，再对来源组取平均：
\[\mathrm{T/B}=\frac{1}{G}\sum_{g=1}^G\frac{T_g}{B_g}.\]

\paragraph{Cmd $\up$：命令完整率。}
公式A只有一个命令\code{\frac}。AST-BPE和软AST-BPE均未在命令内部切开，Cmd为$1/1=100\%$；Qwen2.5将反斜杠与\code{frac}分开，Cmd为$0/1=0\%$。汇总时用完整命令数除以命令总数。

\paragraph{Num $\up$：数字完整率。}
公式B含数字1和12。AST-BPE保留两个数字，Num为$2/2=100\%$；软AST-BPE在12内部切开，Num为$1/2=50\%$。汇总时用内部未切开的数字词素数除以数字词素总数。

\paragraph{Cross $\dn$：语法部分跨越率。}
按KaTeX语法跨度判断token是否部分越过节点。公式A的软AST-BPE将\code{\frac}与左花括号合并，5个token中有1个跨越分子组起点：
\[\mathrm{Cross}_A=1/5\times100\%=20\%.\]
AST-BPE保留分子组，Cross为$0/5=0\%$。实验表用跨越token数除以全部token数，每个token最多计一次。

\paragraph{Used：词表使用率。}
统计输入中出现过的不同ID，除以词表大小：
\[\mathrm{Used}=\frac{\text{已使用的不同ID数}}{\text{词表ID总数}}\times100\%.\]
公式B在软AST-BPE下使用10个不同ID，因此单式Used为$10/1024\times100\%=0.9766\%$。该组完整测试集使用696个ID，Used为$696/1024\times100\%=67.97\%$。
\paragraph{速度 $\up$：完整编码吞吐。}
用原始输入MiB除以完整编码时间：
\[\mathrm{Speed}=\frac{\text{输入字节数}/2^{20}}{\text{编码秒数}}.\]
同组公式A、B所在的完整测试集有659,113字节；软AST-BPE七轮编码的中位耗时为0.02121125秒：
\[\mathrm{Speed}=\frac{659113/2^{20}}{0.02121125}\approx29.63\ \mathrm{MiB/s}.\]
计时包含原子映射、解析和分割，模型加载与指标计算在计时之外。
